\documentclass{article}
\usepackage{amsmath}
\usepackage{graphicx}
\usepackage{fancyhdr}
\usepackage{array}
\usepackage{booktabs}
\usepackage[sorting=none]{biblatex}
\usepackage[margin=1in]{geometry}
\usepackage{listings}


\lstdefinestyle{mystyle}{
    backgroundcolor=\color{backcolour},   
    commentstyle=\color{codegreen},
    keywordstyle=\color{magenta},
    numberstyle=\tiny\color{codegray}, 
    stringstyle=\color{codepurple},
    basicstyle=\ttfamily\footnotesize,
    breakatwhitespace=false,         
    breaklines=true,                 
    captionpos=b,                    
    keepspaces=true,                 
    numbers=left,                    
    numbersep=5pt,                  
    showspaces=false,                
    showstringspaces=false,
    showtabs=false,                  
    tabsize=2,
    frame=single
}

\lstset{style=mystyle}

\lstset{
    language=Python,                    % تنظیم زبان پایتون
    backgroundcolor=\color{white},      % رنگ پس‌زمینه سفید
    basicstyle=\ttfamily\footnotesize,  % استفاده از فونت monospaced
    keywordstyle=\color{blue},           % رنگ کلمات کلیدی
    commentstyle=\color{green},         % رنگ کامنت‌ها
    stringstyle=\color{red},            % رنگ رشته‌ها
    numbers=left,                       % شماره‌گذاری خطوط
    numberstyle=\tiny\color{gray},      % استایل شماره‌های خطوط
    stepnumber=1,                       % شماره‌گذاری هر خط
    numbersep=5pt,                      % فاصله شماره‌گذاری از کد
    showspaces=false,                   % نمایش فضاها
    showstringspaces=false,             % نمایش فضاها در رشته‌ها
    showtabs=false,                     % نمایش تب‌ها
    frame=single,                       % نمایش قاب دور کد
    rulecolor=\color{black},            % رنگ قاب دور کد
    tabsize=2,                          % اندازه تب‌ها
    captionpos=b,                       % عنوان در پایین قرار می‌گیرد
    breaklines=true,                    % تقسیم خط‌های طولانی
    breakatwhitespace=true              % تقسیم در فضای خالی
}
\usepackage[hidelinks]{hyperref}
\usepackage{subfigure}
\hypersetup{
    colorlinks=true,
    linkcolor=teal,
    filecolor=magenta,      
    urlcolor=teal,
    citecolor = teal
    }
\usepackage{xcolor}
\usepackage{xepersian}
\usepackage{fontspec}
\setlength\headheight{28pt} 
\addbibresource{bibliography.bib}
\settextfont[Scale=1.2]{IRLotus.TTF}
\setlatintextfont[Scale=1]{Times New Roman}
\renewcommand{\baselinestretch}{1.5}
\pagestyle{fancy}
\fancyhf{}
\rhead{
\includegraphics[width=1cm]{FaultD.png}مینی پروژه چهارم درس یادگیری ماشین}
\lhead{\thepage}
\rfoot{امیر جهانگرد تکالو}
\lfoot{علیرضا امیری}
\renewcommand{\headrulewidth}{1pt}
\renewcommand{\footrulewidth}{1pt}
\AtBeginDocument{
	\def\chapterautorefname{فصل}%
	\def\sectionautorefname{پاسخ سوال}%
	\def\subsectionautorefname{بخش}%
	\def\subsubsectionautorefname{بخش}%
	\def\equationautorefname{رابطهٔ}%
    \def\lstlistingautorefname{برنامۀ}%
}
\renewcommand{\lstlistingname}{Code}
\begin{document}

\input{titlepage}


\tableofcontents
\newpage


\section{پرسش 1}


\href{https://colab.research.google.com/drive/1oItLI0LIUcmaV6vsL-HBazU3BFLg5jtR?usp=sharing}{\textcolor{blue}{لینک گوگل کولب}}


در ابتدا کتابخانه های مورد نظر را import می کنیم. در ادامه یک کلاس تعریف میکنیم که نورون ما هست. این نورون ورودی های x و y را می گیرد و به کمک وزن ها (Weights) یک مقدار را محاسبه می کند و این مقدار را با مقدار threshold مقایسه می کند و یک خروجی 0 یا 1 به ما می دهد.

\begin{lstlisting}[language=Python]
    class McCulloch_Pitts_neuron():
        def __init__(self, weights, threshold):
            self.weights = weights
            self.threshold = threshold
    
        def model(self, x):
            if self.weights @ x >= self.threshold:
                return 1
            else:
                return 0
\end{lstlisting}

در ادامه تابعی برای بررسی داخل مثلث بودن مینویسیم. بدین صورت که مثلث دارای سه راس است برای اینکه نقطه داخل مثلث باشد باید سه شرط را بررسی کنیم:

\begin{lstlisting}[language=Python]
    def Area(x, y):
        #  y >= 0 → w = [0, 1], threshold = 0
        neur1 = McCulloch_Pitts_neuron([0, 1], 0)
    
        #  y <= 2x - 2 → 2x - y - 2 >= 0 → w = [2, -1], threshold = 2
        neur2 = McCulloch_Pitts_neuron([2, -1], 2)
    
        #  y <= -2x + 6 → -2x - y + 6 >= 0 → w = [-2, -1], threshold = -6
        neur3 = McCulloch_Pitts_neuron([-2, -1], -6)
    
        # Output of each condition
        z1 = neur1.model(np.array([x, y]))
        z2 = neur2.model(np.array([x, y]))
        z3 = neur3.model(np.array([x, y]))
    
        # Final neuron (AND Neuron)
        neur_final = McCulloch_Pitts_neuron([1, 1, 1], 3)
        z_final = neur_final.model(np.array([z1, z2, z3]))
    
        return [z_final]
\end{lstlisting}

برای هر شرط سه نورون تعریف می کنیم:

\begin{lstlisting}[language=Python]
        # Output of each condition
        z1 = neur1.model(np.array([x, y]))
        z2 = neur2.model(np.array([x, y]))
        z3 = neur3.model(np.array([x, y]))
\end{lstlisting}

نرونها را روی نقطه اجرا می کنیم و خروجی میگیریم (0) یا (1). در آخر اگر همه شرطها با هم درست بودن یک نرون دیگه میخوانیم که بررسی کنه همه اینها با هم 1 هستند یا خیر:

\begin{lstlisting}[language=Python]
        # Final neuron (AND Neuron)
        neur_final = McCulloch_Pitts_neuron([1, 1, 1], 3)
        z_final = neur_final.model(np.array([z1, z2, z3]))
    
        return [z_final]
\end{lstlisting}

در ادامه 2000 نقطه در بازه مورد نظر تولید می کنیم:


\begin{lstlisting}[language=Python]
    num_points = 2000
    x_values = np.random.uniform(0, 4, num_points)
    y_values = np.random.uniform(-1, 3, num_points)
\end{lstlisting}

سپس هر نقطه را میدهیم به تابع Area و بررسی می کنیم داخل مثلث هست یا نه. اگر هست، می گذاریم داخل لیست نقاط سبز در غیر این صورت داخل لیست نقاط قرمز. 


\begin{lstlisting}[language=Python]

    # Lists of points based on neuron output
    
    red_points = []
    green_points = []
    
    # Checking any point with the Area function
    for i in range(num_points):
        z = Area(x_values[i], y_values[i])
        if z == [1]:
            green_points.append((x_values[i], y_values[i]))
        else:
            red_points.append((x_values[i], y_values[i]))
    
    # Separating x and y values ​​for plotting
    red_x, red_y = zip(*red_points)
    green_x, green_y = zip(*green_points)
    
\end{lstlisting}

:در نهایت نمودار خروجی را رسم می کنیم

\begin{lstlisting}[language=Python]
    # Plotting
    plt.figure(figsize=(8, 6))
    plt.scatter(red_x, red_y, color='red', label='z = 0')
    plt.scatter(green_x, green_y, color='green', label='z = 1')
    plt.xlabel('X-axis')
    plt.ylabel('Y-axis')
    plt.title('McCulloch-Pitts Neuron Outputs')
    
    # Drawing the sides of triangle ABC
    A, B, C = (2, 2), (3, 0), (1, 0)
    triangle = plt.Polygon([A, B, C], fill=False, edgecolor='black', linestyle='--', linewidth=1.5)
    plt.gca().add_patch(triangle)
    
    # Axis range
    plt.xlim(0, 4)
    plt.ylim(-1, 3)
    plt.grid(True)
    plt.legend(loc='upper right')
    plt.savefig('triangle_output.png', bbox_inches='tight')
    plt.show()
    \end{lstlisting}

مشاهده خروجی:

\begin{figure}[h!]
    \centering
    \includegraphics[width=0.75\linewidth]{1.png}
    \caption{}
    \label{1}
\end{figure}

\newpage

\section{پرسش 2}

\subsection{2.1.1}

طبق چیزی که در این مقاله توضیح داده شده است، داده‌ها از طریق سرویس OpenWeather API جمع آوری شده است. نویسندگان این مقاله برای جمع آوری داده های مربوط به 4 منطقه مختلف از کشور موسیس استفاده کردند: Quatre Bornes, Curepipe, Vacoas, Moka. انجام شد و برای هر منطقه در هر 15 دقیقه یک نمونه داده ثبت شد که در نهایت برای هر منطقه در حدود 2976 نمونه ثبت شد. داده های جمع آوری شده شامل 6 ویژگی مهم آب و هوایی هستند:

\begin{itemize}
\item دما
\item سرعت باد
\item جهت باد
\item فشار هوا
\item رطوبت
\item میزان ابری بودن
\end{itemize}

برای ذخیره سازی این اطلاعات از دو مسیر استفاده می‌شود:

\begin{itemize}
\item ذخیره محلی در پایگاه داده ی MySQL
\item ذخیره در پایگاه داده ابری IBM Cloudant
\end{itemize}

دستگاه‌هایی که در جمع آوری این داده ها سهم بودند شامل اپلیکیکیشن‌های موبایل و دستکتاب بودند که به صورت مستقیم یا از طریق سرویس محلی و ابری، داده‌ها به پایگاه داده‌ها منتقل می‌گردند. هدف از جمع آوری این نوع داده‌ها استفاده از آن‌ها برای آموزش مدل‌های ماشین لرنینگ به صورت collaborative بود به عبارت دیگر استفاده همزمان از داده‌های چند منطقه برای بهبود دقت پیش‌بینی وضعیت آب و هوا در یک منطقه خاص. این ساختار باعث شد دیتاست مقاله از نظر دقت و هم از نظر پوشش مکانی، شرایط مناسبی برای آموزش مدل‌ها داشته باشد و بتوان روی آن تحلیل‌های خوبی انجام داد.

\subsubsection{2.1.2}


در این مرحله، فایل csv \texttt{weather\_prediction\_dataset\.csv} از گوگل درایو وارد محیط گوگل کولب شد. چون فایل اولیه به صورت zip بود پس با استفاده از تابع \texttt{Zipfile} فایل استخراج و سپس مسیر دقیق فایل csv مشخص شد. در ادامه با استفاده از کتابخانه \texttt{pandas} فایل خوانده شد و با دستور \texttt{df\.head()} ردیف‌های اول مشاهده شد. سپس با استفاده از دستور \texttt{df\.columns} فهرست تمام ستون‌های دیتاست را مشاهده کردیم و مشاهده شد که دیتاست شامل اطلاعات مربوط به شهرهای مختلفی مانند \texttt{BASEL\, STOCKHOLM\, TOURS} می‌باشد. بنابراین برای ادامه پروژه تصمیم گرفته شد که فقط از داده‌های مربوط به چند شهر منتخب مثلا از فرانسه استفاده شود و بقیه شهرها در مراحل بعدی فیلتر شوند.

\subsubsection{2.1.3}

در این مرحله فایل داده ها بررسی می شوند که شامل 3654 ردیف داده مربوط به بازه زمانی از 1 ژانویه 2000 تا 1 ژانویه 2010 میباشد. در مقاله مرجع برای یادگیری مدلهای ماشین لرنینگ از پیش پردازش نرمال سازی ویژگیها برای کاهش اختلاف مقیاس بین پارامترهای مختلف مانند دما و فشار از \lr{sliding window} برای استخراج ورودی های متوالی از روزهای گذشته و ساخت مجموعه داده آموزشی مناسب و تنظیم خروجی مدل بر اساس پیش بینی مقدار دمای میانگین روز آینده استفاده شده است.

\subsubsection{2.1.4}

در این قسمت ابتدا ستون \lr{date} را که به صورت عددی بود با استفاده از تابع \lr{to\_datetime()} به فرمت تاریخ تبدیل می‌کنیم تا امکان فیلتر کردن سال‌ها را داشته باشیم. در ادامه با استفاده از \lr{df['Date'].dt.year < 2000} داده‌هایی که قبل از سال 2009 هستند را به عنوان داده‌های \lr{train} جدا کرده، و داده‌های مربوط به سال 2009 ( را نیز به عنوان داده تست استخراج می‌کنیم و در آخر با استفاده از \lr{shape} اندازه هر دو مجموعه چاپ می‌شود.

برای پیش‌بینی مقدار دمای میانگین یک روز، از اطلاعات چند روز قبل مثلاً با اندازه ۵ روز استفاده می‌کنم یعنی با استفاده از روش 
\lr{slide window}
داده‌های \lr{train} و \lr{test} به شکل ورودی مناسب برای مدل ماشین لرنینگ تبدیل می‌شوند. در ابتدا ویژگی‌های مربوط به شهر \lr{TOURS} شامل 
\lr{temp\_min}، \lr{temp\_mean}، \lr{temp\_max} 
انتخاب می‌شوند. سپس با استفاده از یک حالت، برای هر روز ( از روز ششم به بعد ) اطلاعات ۵ روز قبلی به عنوان ورودی مدل و دمای میانگین روز ششم به عنوان خروجی در نظر گرفته شد که برای داده‌های \lr{train} و \lr{test} به صورت جداگانه انجام می‌شود.

\subsection{2.2}


\subsubsection{2.2.1}

\textbf{مفهوم \lr{Collaborative Machine Learning}:}

یادگیری ماشین مشترک به یک روش پیش‌بینی اشاره دارد که در آن داده‌های جمع‌آوری‌شده از چندین مکان یا منبع مختلف به‌طور همزمان برای آموزش مدل‌ها استفاده می‌شود. این مدل‌ها می‌توانند برای پیش‌بینی متغیرهایی مانند دما، رطوبت و فشار جو در مناطق مختلف استفاده شوند.

در روش‌های سنتی یادگیری ماشین، هر مدل معمولاً برای یک منطقه خاص آموزش داده می‌شود و داده‌های مربوط به همان منطقه برای پیش‌بینی استفاده می‌شود. اما در یادگیری ماشین مشترک، داده‌ها از چندین منطقه مختلف جمع‌آوری می‌شوند و مدل‌ها با استفاده از این داده‌های ترکیبی آموزش داده می‌شوند. این رویکرد می‌تواند به مدل‌ها کمک کند که پیش‌بینی دقیق‌تری برای مناطق خاص انجام دهند، به‌ویژه زمانی که داده‌ها از مناطق مختلف همبستگی دارند.

\textbf{ویژگی‌های یادگیری ماشین مشترک:}

یادگیری ماشین مشترک دارای ویژگی‌های زیر است:
\begin{itemize}
    \item \textbf{ترکیب داده‌ها از چندین مکان}: داده‌ها از چندین منطقه یا مکان مختلف برای آموزش مدل‌ها استفاده می‌شود.
    \item \textbf{مدل‌های مشترک}: به‌جای آموزش مدل‌های مجزا برای هر مکان، از یک مدل مشترک استفاده می‌شود که به‌طور همزمان از داده‌های مختلف بهره می‌برد.
    \item \textbf{بهبود دقت پیش‌بینی}: این رویکرد می‌تواند به بهبود دقت پیش‌بینی‌ها کمک کند، به‌ویژه در شرایطی که داده‌ها در مکان‌های مختلف با یکدیگر همبستگی دارند.
\end{itemize}

\textbf{استفاده از \lr{Collaborative Machine Learning} در مقاله:}

در مقاله‌ای که مورد بررسی قرار گرفته، **یادگیری ماشین مشترک** برای پیش‌بینی وضعیت آب و هوا در مناطق مختلف استفاده شده است. داده‌های جمع‌آوری‌شده از چندین مکان مختلف، مانند **Basel**، **Stockholm** و **Tours**، برای آموزش مدل‌ها استفاده شده‌اند. این مدل‌ها به‌طور خاص برای پیش‌بینی پارامترهای مختلف مانند دما، رطوبت، فشار و تابش خورشیدی استفاده می‌شوند.

\textbf{استفاده از داده‌های مشترک:}

در این مقاله، داده‌های جمع‌آوری‌شده از مکان‌های مختلف (شهرهای مختلف) به‌طور همزمان به مدل‌های یادگیری ماشین وارد می‌شوند. این داده‌ها شامل پارامترهای مختلف آب و هوا مانند دما، رطوبت و فشار هستند. این داده‌ها سپس برای پیش‌بینی وضعیت آب و هوا در یک منطقه خاص مورد استفاده قرار می‌گیرند.

این روش می‌تواند به پیش‌بینی دقیق‌تری منجر شود زیرا اطلاعات بیشتری از مناطق مختلف در دسترس است و مدل‌ها می‌توانند از این اطلاعات برای آموزش استفاده کنند. مقاله نشان می‌دهد که پیش‌بینی‌هایی که از داده‌های مشترک استفاده می‌کنند، دارای دقت بالاتری نسبت به پیش‌بینی‌هایی هستند که فقط از داده‌های یک منطقه خاص استفاده می‌کنند.

\textbf{مدل‌های یادگیری ماشین استفاده‌شده:}

در این مقاله، از پنج مدل مختلف یادگیری ماشین برای پیش‌بینی وضعیت آب و هوا استفاده شده است. این مدل‌ها عبارتند از:
\begin{itemize}
    \item \textbf{رگرسیون خطی چندگانه (MLR)}
    \item \textbf{رگرسیون چندجمله‌ای (MPR)}
    \item \textbf{\lr{K-Nearest Neighbors (KNN)}}
    \item \textbf{پرسپترون چند لایه (MLP)}
    \item \textbf{شبکه عصبی کانولوشنی (CNN)}
\end{itemize}
این مدل‌ها به‌طور خاص برای پیش‌بینی ویژگی‌های مختلف آب و هوا آموزش داده شدند. مقاله نشان داد که استفاده از داده‌های مشترک برای آموزش این مدل‌ها باعث بهبود قابل توجهی در دقت پیش‌بینی‌ها شد.

\textbf{نتایج استفاده از \lr{Collaborative Machine Learning}:}

در این مقاله، نتایج نشان داد که پیش‌بینی‌های مشترک با استفاده از داده‌های چندین مکان به‌طور قابل توجهی دقت پیش‌بینی‌ها را بهبود بخشیده است. در آزمایش‌های انجام‌شده، مدل‌هایی مانند \textbf{رگرسیون چندجمله‌ای (MPR)} و \textbf{رگرسیون خطی چندگانه (MLR)} که از داده‌های مشترک استفاده کرده بودند، عملکرد بهتری نسبت به مدل‌های پیش‌بینی غیر مشترک داشتند. مقاله همچنین نشان داد که پیش‌بینی‌های مشترک می‌توانند میانگین خطای پیش‌بینی را تا حدود 5\% کاهش دهند.

\textbf{استفاده از پیش‌بینی‌های مشترک برای بهبود دقت:}

نتایج این مقاله نشان داد که استفاده از پیش‌بینی‌های مشترک به‌ویژه در مناطقی که داده‌های آن‌ها به‌طور طبیعی همبسته هستند، می‌تواند دقت پیش‌بینی‌ها را بهبود بخشد. این روش از طریق ترکیب داده‌های مختلف و استفاده از یک مدل مشترک به‌جای آموزش مدل‌های جداگانه برای هر منطقه، به دقت بالاتر در پیش‌بینی‌ها کمک می‌کند.

\subsection{2.3}

در این قسمت ما ورودی های مدل را انتخاب کرده به این صورت که رطوبت، فشار هوا و تابش خورشید را از دیتافریم df برداشته و
به عنوان ورودی به مدل داده ایم. در خط بعد مقدار دمای ماکزیمم را به عنوان خروجی مدل تعریف میکنیم. در ادامه کد، دادهها را
به 0 بخش تقسیم کرده به این گونه که 02 \% دادهها برای تست کنار گذاشته میشوند و 02 \% باقیمانده برای آموزش مدل استفاده
میشوند.


در این بخش، یک شبکه عصبی ساده با استفاده از کتابخانه keras طراحی شده است.ساختار این شبکه شامل یک لایه پنهان با 12
نورون است که از تابع فعالساز ReLU برای افزایش توان مدل در یادگیری روابط غیرخطی استفاده میکند. در نهایت، لایه خروجی
تنها یک نورون دارد. در ابتدا کتابخانه های keras برای ساخت مدل فراخوانی میشود. در لیست \lr{learning\_rates} سه مقدار مختلف
نرخ یادگیری وجود دارد که برای مقایسه تاثیر آنها در یادگیری مدل استفاده میشود. مدل با استفاده از الگوریتم بهینهسازی SGD و
تابع هزینه MSE کامپایل میشود. مدل به صورت 20 epoch آموزش داده میشود و پس از آن مقدار Loss برای دادههای train و
test محاسبه و ذخیره میشود.


در لیست اول، سه مقدار مختلف برای نرخ یادگیری ذخیره شده که در آنها 1 مقدار خیلی زیادی است و ممکن است باعث بیثباتی
مدل شود. \lr{0.001} مقدار مناسب و متعادل است که معمولا بهترین انتخاب است و سوم \lr{0.00000001} مقدار بسیار کمی است که مدل
خیلی کند یا اصلا قادر به یادگیری نیست. در لیست دوم مقدار خطای train برای هر نرخ یادگیری نگه داشته میشود که زمانیکه
\lr{learning rate = 0.001} باشد مقدار \lr{train loss = 4.6} میشود که بهترین انتخاب است. در خط آخر نیز مقدار خطای test برای
همان نرخها ذخیره میشود دکه بهترین انتخاب وقتی \lr{learning rate =0.001} باشد، \lr{test loss = 5.0} میشود.

\begin{figure}[h!]
    \centering
    \includegraphics[width=0.75\linewidth]{2.png}
    \caption{نمودار میزان خطا بر اساس نرخ یادگیری}
    \label{2}
\end{figure}

\begin{figure}[h!]
    \centering
    \includegraphics[width=0.75\linewidth]{3.png}
    \caption{نمودار دقیقتر میزان خطا بر اساس نرخ یادگیری}
    \label{3}
\end{figure}

\newpage

\subsection{2.4}

\subsubsection{2.4.1}

\begin{figure}[h!]
    \centering
    \includegraphics[width=0.75\linewidth]{4.png}
    \caption{نمودار دقیقتر میزان خطا بر اساس تمامی حالات نرخ یادگیری}
    \label{4}
\end{figure}

\begin{figure}[h!]
    \centering
    \includegraphics[width=0.75\linewidth]{5.png}
    \caption{میزان خطای آموزش و اعتبار سنجی با معیار MSE}
    \label{5}
\end{figure}

\newpage

\subsubsection{2.4.2}

به وضوح می توان مشاهده کرد که با افزایش نرخ یادگیری میزان خطای ما کاهش می یابد و نزدیک صفر می شود که این یک نکته مثبت است.



\subsection{2.5}

\textbf{تغییرات شبکه:}

\begin{itemize}
    \item \textbf{افزایش تعداد لایه‌ها}: از ۲ لایه به ۵ لایه افزایش یافت
    \item لایه‌های پنهان: ۶۴، ۱۲۸، ۶۴ و ۳۲ نورون
    \item لایه خروجی: ۱ نورون (برای رگرسیون)
    \item \textbf{الگوی معماری}: ایجاد ساختار ساعت‌شنی (بزرگ به کوچک)
    
\end{itemize}


\begin{itemize}
    \item کاهش تدریجی تعداد نورون‌ها در لایه‌های عمیق‌تر
\end{itemize}

\begin{itemize}
\item اضافه شدن لایه‌های Dropout با نرخ‌های مختلف (۰.۳ و ۰.۴)
\item جلوگیری از overfitting با غیرفعال کردن تصادفی نورون‌ها
\end{itemize}

\textbf{تنظیم L2}:

\begin{itemize}
    \item اعمال جریمه بر وزن‌های بزرگ
    \item ضریب تنظیم: ۰.۰۰۱
\end{itemize}

\textbf{\lr{Batch Normalization}}:

\begin{itemize}
    \item نرمال‌سازی خروجی هر لایه
    \item پایداری بیشتر فرآیند آموزش
    \item امکان استفاده از نرخ یادگیری بالاتر
\end{itemize}

\textbf{بهینه‌ساز Adam}:

\begin{itemize}
\item جایگزین SGD شده است
\item نرخ یادگیری تطبیقی برای هر پارامتر
\item نرخ یادگیری اولیه: ۰.۰۰۱
\end{itemize}

\textbf{افزایش تعداد اپاک}:
\begin{itemize}
    \item از ۲۰ به ۱۰۰ اپاک افزایش یافت
    \item وجود EarlyStopping باعث می‌شود آموزش زودتر متوقف شود اگر بهبودی حاصل نشود
\end{itemize}

\textbf{تعیین اندازه بچ}:
\begin{itemize}
    \item اندازه بچ ۳۲ برای به روزرسانی وزن‌ها
\end{itemize}

\textbf{EarlyStopping}:

\begin{itemize}
\item نظارت بر loss اعتبارسنجی
\item توقف آموزش پس از ۱۰ اپاک بدون بهبود (patience=10)
\item بازگرداندن بهترین وزن‌ها
\end{itemize}

\textbf{ReduceLROnPlateau}:
\begin{itemize}
    \item کاهش خودکار نرخ یادگیری هنگام توقف بهبود
    \item ضریب کاهش: ۰.۲
    \item حداقل نرخ یادگیری: ۱e-6
\end{itemize}

\begin{figure}[h!]
    \centering
    \includegraphics[width=1\linewidth]{6.png}
    \caption{نمودار دقیقتر میزان خطا برای مدل عمیق تر بر اساس تمامی حالات نرخ یادگیری و مقایسه مدل های پیشبینی شده با مقدار اصلی}
    \label{6}
\end{figure}

\begin{figure}[h!]
    \centering
    \includegraphics[width=0.75\linewidth]{7.png}
    \caption{میزان خطای آموزش و اعتبار سنجی برای مدل عمیق تر با معیار MSE}
    \label{7}
\end{figure}

\newpage

مشاهده می کنیم که خطای ما خیلی زودتر از حالت قبل به صفر میل می شود. در ضمن میزان خطای MSE در مدل عمیقتر خیلی کمتر از مدل عادی است.


\subsection{2.6}


\begin{figure}[h!]
    \centering
    \includegraphics[width=0.75\linewidth]{8.png}
    \caption{نمودار تغییرات وزن ها بر هر EPOCHS}
    \label{8}
\end{figure}

\newpage

اگر وزن ها به مرور به یک مقدار ثابت نزدیک شوند می توان گفت که مدل در حال یادگیری پایدار است ولی اگر وزن ها نوسان زیادی داشته باشند یا هیچ تغییر خاصی نکنند ممکن است مشکلی در یادگیری وجود داشته باشد. مثلا نرخ یادگیری نادرست یا مدل ضعیف. پس تاثیر انتخاب نرخ یادگیری درست بسیار مهم است

همچنین ما در قسمت قبل دیدیم که تاثیر مواردی مثل Dropout و Regularization و \lr{Batch Normalization} چقدر در میزان خطای خروجی موثر است.

\section{پرسش 3}

\href{https://drive.google.com/file/d/1rySLMOsEJoQS4XSUzd8LL4EnznuVOkk2/view?usp=sharing}{\textcolor{blue}{لینک گوگل کولب}}

در این پرسش، به طرح یک مدل یادگیری تقویتی در فضای مسئله ی wumpos خواهیم پرداخت. در این مسئله، عامل تعریف شده می تواند در چهار جهت راست، چپ، بالا و پایین حرکت کند و بر اساس موقعیت بعدی، میزانی پاداش مثبت و یا منفی دریافت کند. 
در صورتی که این عامل به یک بلوک خالی برود امتیازی دریافت نمی کند، در صورتی که در گودال (pit) بیافتد و یا توسط wumpus خورده شود، 1000 امتیاز منفی، در صورت رسیدن به طلا 100 امتیاز مثبت و به ازای هر حرکت، 1 امتیاز دریافت می کند. 
در گام اول، دنیای شبیه سازی به صورت نمایش داده شده در \autoref{fig:9} طراحی می شود. 
\begin{figure}
	\centering
	\includegraphics[width=0.7\linewidth]{9}
	\caption{Wompus World}
	\label{fig:9}
\end{figure}
\pagebreak
\subsection{آموزش مدل بر اساس الگوریتم Q-Learning}
پس از تعریف فضای مدل و اکشن های آن، به تعریف مدل یادگیری تقویتی خواهیم پرداخت. در این بخش، ساختار مدل بر اساس Q-Leaning با پارامتر های زیر تعریف می شود.
\begin{table}[h!]
	\centering
	\caption{پارامترهای استفاده شده در الگوریتم Q-Learning}
	\begin{tabular}{|c|c|c|}
		\hline
		\textbf{نام پارامتر} & \textbf{مقدار پیش‌فرض} & \textbf{توضیح} \\
		\hline
		\lr{env} & ورودی از کلاس محیط & محیط آموزش عامل (مثلاً \lr{WumpusWorld}) \\
		\hline
		\lr{alpha} & $0.1$ & نرخ یادگیری (یادگیری از تجربه‌های جدید) \\
		\hline
		\lr{gamma} & $0.9$ & ضریب تخفیف برای پاداش‌های آینده \\
		\hline
		\lr{epsilon} & $1.0$ & نرخ اکتشاف؛ احتمال انتخاب تصادفی عمل \\
		\hline
		\lr{epsilon\_decay} & $0.995$ & ضریب کاهش تدریجی \lr{epsilon} در هر اپیزود \\
		\hline
		\lr{min\_epsilon} & $0.01$ & حداقل مقدار \lr{epsilon} برای حفظ اکتشاف جزئی \\
		\hline
		\lr{q\_table} & \lr{np.zeros((4,4,4))} & جدول Q با حالت‌ها و اعمال ممکن \\
		\hline
		\lr{actions} & لیست اعمال محیط & مثل \lr{['up', 'down', 'left', 'right']} \\
		\hline
		\lr{rewards} & لیست خالی & ذخیره مجموع پاداش اپیزودها برای تحلیل \\
		\hline
	\end{tabular}
\end{table}

\subsection*{تابعی برای نمایش Policy}

برای بررسی عملکرد عامل یادگیری تقویتی پس از اتمام فرایند آموزش، تابعی به نام \lr{show\_policy} پیاده‌سازی شده است که قادر است سیاست یادگرفته‌شده را به‌صورت تصویری روی محیط چهار در چهار نشان دهد. این نمایش شامل نمادهایی از جهت حرکت (پیکان‌ها) و مسیر بهینه عامل از نقطه شروع تا رسیدن به وضعیت هدف است.

در ابتدا، ماتریسی به‌نام \lr{policy\_grid} ساخته می‌شود که در آن برای هر خانه، بهترین عمل بر اساس بیشترین مقدار Q ثبت می‌گردد. اعمال به‌صورت پیکان‌های جهت‌دار مانند نمایش داده می‌شوند. در خانه‌هایی که شامل طلا، گودال یا Wumpus باشند، به‌ترتیب ایموجی‌های مربوطه قرار داده می‌شود تا تشخیص وضعیت‌ها آسان‌تر گردد.

در صورتی که گزینه \lr{show\_best\_path} فعال باشد، مسیر بهینه‌ای که عامل در صورت اجرای greedy (انتخاب همیشگی بهترین عمل) از موقعیت شروع طی می‌کند نیز محاسبه و روی محیط نمایش داده می‌شود. این مسیر با استفاده از نقاط آبی رنگ مشخص می‌گردد.

هدف از این تابع، فراهم‌کردن یک ابزار شهودی برای بررسی کیفیت سیاست یادگرفته‌شده توسط عامل است. با استفاده از آن می‌توان مشاهده کرد که آیا عامل موفق شده است از موانع عبور کند، از خطرات اجتناب نماید و در نهایت به‌طور مؤثر به طلا برسد یا خیر.

با در اختیار داشتن فضای شبیه سازی، مدل یادگیری و همچنین تابعی برای نمایش نتیجه ی بهترین مدل، شروع به آموزش مدل می کنیم. 
در \autoref{fig:10} نتیجه ی این مدل به همراه بهترین مسیر به دست آمده پس از 500 تکرار نمایش داده شده است.
\begin{figure}
	\centering
	\includegraphics[width=0.7\linewidth]{10}
	\caption{نتیجه مدل Q-Learning}
	\label{fig:10}
\end{figure}

\pagebreak

همچنین، میزان امتیاز هر مرتبه آزمایش در فرایند یادگیری در \autoref{fig:11} رسم شده است.
\begin{figure}
	\centering
	\includegraphics[width=0.7\linewidth]{11}
	\caption{امتیازات مدل در فرایند یادگیری}
	\label{fig:11}
\end{figure}

\pagebreak

\subsection{آموزش مدل DQN}
در گام بعد، به آموزش یک مدل DQN خواهیم پرداخت. در این روش، ساختار مدل به صورت زیر طراحی شده است:\begin{figure}
	\centering
	\includegraphics[width=0.7\linewidth]{12}
	\caption{ساختار مدل DQN}
	\label{fig:12}
\end{figure}
\pagebreak

در این ساختار، ورودی 16 تایی از موقعیت های ممکن به صورت $one hot coded$ به مدل داده می شود و پس از دو لایه پنهان متشکل از نورون های 64 تایی، خروجی به صورت هشت تایی متشکل از اکشن های حرکت در راستا های چپ، راست، بالا و پایین  و نیز شلیک در این راستا ها ثبت می شود. 
در این مدل، شبکه عصبی جایگزین تابع Q می شود و به جای یادگیری مقادیر موجود در آن تابع، وزن های شبکه آموزش می بینند.
مجموع امتیازات این مدل را در \autoref{fig:13} می توانیم مشاهده کنیم.
\begin{figure}
	\centering
	\includegraphics[width=1\linewidth]{13}
	\caption{}
	\label{fig:13}
\end{figure}
\pagebreak


\subsection{عملکرد Policy}

در این بخش، عملکرد سیاست یادگرفته‌شده در طول اپیزودهای مختلف مورد بررسی قرار گرفته است. برای این منظور، مجموع امتیازات کسب‌شده توسط عامل در هر اپیزود ثبت شده‌است تا مشخص شود که آیا عامل به مرور زمان قادر به یادگیری حرکت بهینه در محیط هست یا خیر.

در ابتدای آموزش، به دلیل بالا بودن نرخ اکتشاف (epsilon)، عامل اغلب حرکات تصادفی انجام می‌دهد و در نتیجه عملکرد ضعیفی دارد؛ زیرا بیشتر حرکات منجر به افتادن در چاله، خورده شدن توسط وامپوس یا عدم دستیابی به طلا می‌شود.

با پیشرفت در آموزش و کاهش تدریجی نرخ اکتشاف، عامل کمتر به صورت تصادفی عمل کرده و بیشتر بر اساس سیاست یادگرفته‌شده تصمیم‌گیری می‌کند. این روند باعث افزایش امتیازها و بهبود عملکرد عامل می‌شود، تا جایی که در اپیزودهای پایانی، عامل قادر است به صورت مؤثر به سمت طلا حرکت کرده و از موانع خطرناک اجتناب کند.

در نتیجه، می‌توان گفت که مدل \lr{DQN} توانسته است سیاست بهینه‌ای را برای حرکت در محیط وامپوس یاد بگیرد و عملکرد قابل قبولی را از خود نشان دهد.

\subsection{تاثیر اپسیلون}
در صورت بالاتر بودن مقدار اپسیلون، میزان exploration مدل افزایش می یابد. به این معنا که در هر حرکت، احتمال انجام دادن حرکتی رندوم که حتی پاداش کمتری نیز دارد بیشتر خواهد بود. این امر منجر می شود که مدل مسیر های متفاوت را شناسایی کرده و بتواند درک بهتری از فضای پیرامون خود داشته باشد. اگرچه این کار منجر به دریافت امتیازات کمتری می شود. با ادامه ی فرایند آموزش، میزان اپسیلون کاهش پیدا کرده و مدل از مسیر های شناخته شده که امتیاز بیشتری داشته اند اسقبال بیشتری می کند و کمتر به جستجوی مسیر های دیگر می پردازد.

\subsection{بررسی کارایی}
مدل Qlearning پس از 800 اپیزود آزمون و خطا توانست به طور دقیق طلا را پیدا کرده و بالاترین امتیاز را به دست آورد. 

در مراحل اولیه‌ی آموزش DQN، به دلیل بالا بودن نرخ اکتشاف، عامل عمدتاً حرکات تصادفی انجام می‌دهد و عملکرد ضعیفی دارد. اما با کاهش تدریجی این نرخ، عامل تجربه‌های قبلی را به‌کار می‌گیرد و به مرور زمان مسیرهای بهینه‌تری را انتخاب می‌کند. در نتیجه، امتیازات کسب‌شده افزایش یافته و عامل قادر به اجتناب از خطرات و رسیدن به طلا می‌شود.

در مقایسه با مدل \lr{Q-Learning}، مدل \lr{DQN} عملکرد سریع‌تر و پایدارتری از خود نشان داده و میانگین امتیازات بالاتری کسب کرده است.


\subsection{ساختار مدل DQN}
توضیحات مربوط به این بخش در قسمت های قبل ذکر شده است.


 













































\end{document}